\section{The same original weighted points compute intrinsic line-section order}

The central curve, ambient scheme, their original stalks and every point in
this statement are unchanged.  This is the actual norm/order dictionary,
not a change of the test curve and not an assumed comparison of slopes.

\begin{theorem}[The actual original line norm has the native intrinsic-order profile]
\label{thm:intrinsic-numerical-gauss-original-line-order-profile}
\lean{EnvelopeProperty.NumericalGaussCurveOrderSlopes.exists_eventual_modelNorm_intrinsicOrder}
\leanok
\uses{spectrum:gauss-valuation,thm:intrinsic-numerical-central-section-original-order-and-norm,thm:intrinsic-numerical-two-frame-leading-line-order,thm:gauss-weighted-existing-native-true-center,lem:actual-generic-native-field-valuations,lem:gauss-native-original-stalk-point-sequence,lem:gauss-native-original-unit-profiles,lem:gauss-surface-native-order-real-log,thm:gauss-surface-actual-fixed-radius-points,def:curve-generic-section-original-order,def:gauss-central-line-original-generic-carrier,cor:central-generic-prescribed-section-same-rational-lift,def:central-generic-original-local-section-value}
Let $Y$ be original integral and proper over $k[t]$, with $k$ a field,
and let $i:C\hookrightarrow Y$ be the original closed immersion of an
integral locally Noetherian curve with dimension at most one.  Fix ONE
original invertible line $L$, ONE actual frame $\delta$ on a neighborhood
$U$ of $i(\eta_C)$, and its SAME original generic sections $\rho,\sigma_C$.
For an original nongeneric $x\in C$, put $p=i(x)$ and
$S=\mathcal O_{Y,p}$.  Assume the actual ordinary local geometry:
$S$ Noetherian, $\pi$ prime, $\pi\nmid z$,
$\mathfrak m_S=(\pi,z)$, and the kernel of the ORIGINAL stalk map
$S\to\mathcal O_{C,x}$ equal to $(\pi)$.  Assume also the original
parameter $t_p\ne0$ and $\pi\mid t_p$.  These are ordinary geometric
hypotheses; no order, residue, norm or point matching is data.

For $0<r<1$, let $q_N$ be the SAME already constructed original native
fixed-Gauss weighted point and let $\epsilon_N$ be its SAME actual positive
native scale.  There is $M$ such that for EVERY $N>0$ with $M<N$,
\[
 -\log\|\rho(q_N)\|_L
   =\epsilon_N\,\operatorname{order}(i^*L,\sigma_C,x).
\]
The norm is evaluated on the original generic-fibre section, using its
actual native point-equality transport.  Its intrinsic central order is
that of the SAME original $\sigma_C$.  In particular $p\in U$ or $p\in W$
is not an assumption, and no norm-one premise is needed for this profile.
\end{theorem}
\begin{proof}
The actual Gauss maximum is bounded by one: each nonzero coefficient
has trivial-base norm one and every radius power is at most one.
Internally apply the actual central coefficient/order and true-center norm
producer to the fixed $\delta$ and original $x$.  It gives the actual
$p$-frame on $V\ni p$, the original overlap $W$ containing $i(\eta_C)$,
a genuine regular unit $r_W$, the actual component-generic stalk unit
$u$, and SAME actual nonzero $f\in K(Y)$, $b\in K(C)$.  It proves
$f=\operatorname{germ}_{\eta_Y}(r_W)=\operatorname{algebraMap}(u)$,
$b=i^\#_{\eta_C}u$, and the intrinsic order of $\sigma_C$ equal to
$C.\operatorname{ord}(b,x)$.

The genuine component-unit leading-order theorem applies to THIS actual
$u$.  Its internally derived original vertical fraction and canonical
residue squares give vertical order zero and eventual actual weighted
order $\exp(C.\operatorname{ord}(b,x))$ on the SAME original stalk fraction
field.  The old native field map is exactly the inverse of the original
stalk/function-field fraction isomorphism, so it sends THIS $f$ to
that actual fraction.  Nonvanishing follows from the already proved
$f\ne0$ and this actual field isomorphism.

The independently proved weighted-center theorem gives the TRUE center
of EVERY $q_N$ equal to $p$; its old native support theorem gives the
original generic point.  Thus the retained norm formula computes
$\|\rho(q_N)\|$ by the original whole-field valuation of $f$.
Use the genuine $W$-regular unit representation of $f$ to apply the old
whole generic-valuation section-evaluation identity and the original
native weighted-point evaluation theorem.  The latter works on EVERY
nonempty original open and does NOT require $p\in W$.  These actual
section/germ triangles identify this value with the old normalized
weighted valuation of the SAME original fraction.

That normalized weighted point uses exactly the native scale
$\epsilon_N$; its actual positive parameter order proves this scale is
positive.  The genuine native-order logarithm formula therefore gives
$\log\|\rho(q_N)\|=-\epsilon_N C.\operatorname{ord}(b,x)$ eventually.
Negate it and substitute the already proved intrinsic-order identity.
This proves the literal norm/order profile without a supplied comparison.
\end{proof}

The final weighted-degree consumer must use the ONE fixed original
$\sigma_C$ and its intrinsic whole-curve Euler-degree formula.  The varying
local functions $f_x$ are not separately subjected to principal balance.
The profile alone claims neither global nonnegative degree nor the final
LocalPositivity theorem.
